\chapter{Subprogram Declarations\label{chap:SubprogramDeclarations}}

The \Tfunc{} keyword declares a subprogram. If a return type is included, it is a \emph{function declaration}.
Otherwise, it is a \emph{procedure declaration}.

\ExampleDef{Subprogram Declarations}
In \listingref{SubprogramDeclarations}, \verb|Add| is a function whereas \verb|IncrementCounter| is a procedure.
\ASLListing{Subprogram declarations}{SubprogramDeclarations}{\definitiontests/SubprogramDeclarations.asl}

\listingref{ParametricFunction} shows an example of a parametric function \verb|f| and its invocations.
\ASLListing{Invoking a parametric function}{ParametricFunction}{\definitiontests/ParametricFunction.asl}

\listingref{ParametricFunction2} shows an example of a parametric function \verb|append| and its invocations.
\ASLListing{Invoking a parametric function}{ParametricFunction2}{\definitiontests/ParametricFunction2.asl}

\ChapterOutline
\begin{itemize}
  \item \secref{Accessors} discusses \emph{accessors};
  \item \FormalRelationsRef{Subprogram Declarations} defines the formal relations for subprogram declarations;
  \item \SyntaxRef{Subprogram Declarations} defines the syntax of subprogram declarations;
  \item \AbstractSyntaxRef{Subprogram Declarations} defines the abstract syntax of subprogram declarations; and
  \item \TypeRulesRef{Subprogram Declarations} defines the type rules for subprogram declarations.
\end{itemize}

%%%%%%%%%%%%%%%%%%%%%%%%%%%%%%%%%%%%%%%%%%%%%%%%%%%%%%%%%%%%%%%%%%%%%%%%%%%%%%%%
\section{Accessors\label{sec:Accessors}}
Special kinds of subprograms are \emph{getters} and \emph{setters}, which are paired to form an \emph{accessor}
declaration. Intuitively, an accessor is used to abstract reading and writing to a global storage datatype.

\listingref{AccessorDeclaration} shows an accessor declaration.
The accessor \texttt{X} is used to read and write underlying storage element \texttt{R}.
For example, this could model a register file consisting of 32 registers, each of 64-bits, the last of which is always zero.

\ASLListing{An accessor declaration}{AccessorDeclaration}{\definitiontests/Accessor.asl}

%%%%%%%%%%%%%%%%%%%%%%%%%%%%%%%%%%%%%%%%%%%%%%%%%%%%%%%%%%%%%%%%%%%%%%%%%%%%%%%%
\FormalRelationsDef{Subprogram Declarations}
\paragraph{Syntax:} Subprogram declarations are grammatically derived from $\Ndecl$ via the subset of productions shown in
  \SyntaxRef{Subprogram Declarations}.

\paragraph{Abstract Syntax:} Subprogram declarations are derived in the abstract syntax from $\decl$
  and generated by $\builddecl$ (see \AbstractSyntaxRef{Subprogram Declarations}).

\paragraph{Typing:} Subprogram declarations are annotated via $\annotateanddeclarefunc$ (see \TypingRuleRef{AnnotateAndDeclareFunc});

\paragraph{Semantics:} The dynamic semantics of subprogram calls is given in \chapref{SubprogramCalls}.

%%%%%%%%%%%%%%%%%%%%%%%%%%%%%%%%%%%%%%%%%%%%%%%%%%%%%%%%%%%%%%%%%%%%%%%%%%%%%%%%
\SyntaxDef{Subprogram Declarations}
%%%%%%%%%%%%%%%%%%%%%%%%%%%%%%%%%%%%%%%%%%%%%%%%%%%%%%%%%%%%%%%%%%%%%%%%%%%%%%%%

\RequirementDef{NoEmptyBody}
Subprograms may not have empty bodies.
To define a subprogram whose body has no effect, a \passstatementterm{} can be used.
For example, \listingref{NoEmptyBody} shows a subprogram declaration whose body has no effect.
\ASLListing{A pass statement can be used to define a subprogram body with no effect}{NoEmptyBody}{\definitiontests/GuideRule.NoEmptyBody.asl}

\begin{flalign*}
\Ndecl  \derives \ & \Npuritykeyword \parsesep \Noverride \parsesep \Tfunc \parsesep \Tidentifier \parsesep \Nparamsopt \parsesep \Nfuncargs &\\
    &\wrappedline \Nreturntype \parsesep \Nrecurselimit \parsesep \Nfuncbody\\
|\ & \Nqualifier \parsesep \Noverride \parsesep \Tfunc \parsesep \Tidentifier \parsesep \Nparamsopt \parsesep \Nfuncargs &\\
    &\wrappedline \Nrecurselimit \parsesep \Nfuncbody &\\
|\ & \Noverride \parsesep \Taccessor \parsesep \Tidentifier \parsesep \Nparamsopt \parsesep \Nfuncargs \parsesep \Tbeq &\\
   & \wrappedline\ \Tidentifier \parsesep \Nasty \parsesep \Naccessorbody &\\
\Naccessorbody \derives \ & \Tbegin \parsesep \Naccessors \parsesep \Tend \parsesep \Tsemicolon&
\end{flalign*}

\begin{flalign*}
\Nrecurselimit   \derives \ & \Trecurselimit \parsesep \Nexpr &\\
|\              & \emptysentence &\\
\Nparamsopt \derives \ & \emptysentence &\\
                   |\ & \Tlbrace \parsesep \ClistZero{\Nopttypedidentifier} \parsesep \Trbrace &\\
\Nopttypedidentifier \derives \ & \Tidentifier \parsesep \option{\Nasty} &\\
\Nfuncargs          \derives \ & \Tlpar \parsesep \ClistZero{\Ntypedidentifier} \parsesep \Trpar &\\
\Nreturntype        \derives \ & \Tarrow \parsesep \Nty &\\
\Nfuncbody          \derives \ & \Tbegin \parsesep \Nstmtlist \parsesep \Tend \parsesep \Tsemicolon &\\
\Nstmtlist \derives \ & \emptysentence \;|\; \Nstmtlist &\\
  \Naccessors \derives \ &
    \Nisreadonly \parsesep \Tgetter \parsesep \Nstmtlist \parsesep \Tend \parsesep \Tsemicolon \parsesep \\ & \wrappedline
    \Tsetter \parsesep \Nstmtlist \parsesep \Tend \parsesep \Tsemicolon &\\
  |\ & \Tsetter \parsesep \Nstmtlist \parsesep \Tend \parsesep \Tsemicolon \parsesep \\ & \wrappedline
      \Nisreadonly \parsesep \Tgetter \parsesep \Nstmtlist \parsesep \Tend \parsesep \Tsemicolon &\\
\Nqualifier \derivesinline\ & \emptysentence \;|\; \Tpure \;|\;\Treadonly \;|\;\Tnoreturn &\\
\Npuritykeyword \derivesinline\ & \emptysentence \;|\; \Tpure \;|\;\Treadonly &\\
\Nisreadonly \derivesinline\ & \emptysentence \;|\; \Treadonly &\\
\Noverride \derivesinline\ & \emptysentence \;|\; \Timpdef \;|\;\Timplementation &
\end{flalign*}

%%%%%%%%%%%%%%%%%%%%%%%%%%%%%%%%%%%%%%%%%%%%%%%%%%%%%%%%%%%%%%%%%%%%%%%%%%%%%%%%
\AbstractSyntaxDef{Subprogram Declarations}
%%%%%%%%%%%%%%%%%%%%%%%%%%%%%%%%%%%%%%%%%%%%%%%%%%%%%%%%%%%%%%%%%%%%%%%%%%%%%%%%
\RenderTypes[remove_hypertargets]{decl_func}


\ASTRuleDef{SubprogramDecl}
This rule defines the cases for $\builddecl$ for subprogram declarations.

We first define $\accessorpair$, which we use in this section as an intermediate representation between the syntax forms of accessors and their corresponding abstract syntax.
In particular, rather than directly building the abstract syntax for accessors, we first build an $\accessorpair$ structure, which we then desugar into abstract syntax.

\hypertarget{ast-accessorpair}{}
\begin{flalign*}
\accessorpair \derives\ &
{
\left\{
  \begin{array}{rcl}
    \accessorpairisreadonly &:& \Bool, \\
    \accessorpairgetter &:& \stmt, \\
    \accessorpairsetter &:& \stmt
\end{array}
\right\}
} &
\end{flalign*}

\begin{mathpar}
\inferrule[func\_decl]{
  \buildqualifier(\vp) \astarrow \astversion{\vp}
}{
  \builddecl\left(
    \overname{\Ndecl\left(
      \begin{array}{l}
    \namednode{\vp}{\Npuritykeyword}, \punnode{\Noverride}, \Tfunc, \Tidentifier(\name), \\
    \wrappedline\ \punnode{\Nparamsopt}, \punnode{\Nfuncargs}, \punnode{\Nreturntype}, \\
    \wrappedline\ \punnode{\Nrecurselimit}, \punnode{\Nfuncbody}
      \end{array}
  \right)}{\vparsednode}\right)
\astarrow \\
{
\overname{
 \left[\DFunc\left(\left\{
  \begin{array}{rcl}
  \FIELDname &:& \name,\\
  \FIELDparameters &:& \astof{\vparamsopt},\\
  \FIELDargs &:& \astof{\vfuncargs},\\
  \funcbody &:& \astof{\vfuncbody},\\
  \FIELDreturntype &:& \some{ \astof{\vreturntype} },\\
  \funcsubprogramtype &:& \STFunction,\\
  \FIELDrecurselimit   &:& \astof{\vrecurselimit}\\
  \funcbuiltin &:& \False\\
  \funcqualifierfield &:& \astversion{\vp}\\
  \funcoverride &:& \astof{\voverride}
 \end{array}
 \right\}\right)\right]}{\vastnode}
}}
\end{mathpar}

\begin{mathpar}
\inferrule[procedure\_decl]{}{
 \builddecl\left(
  \overname{\Ndecl\left(
   \begin{array}{l}
   \punnode{\Nqualifier}, \punnode{\Noverride}, \Tfunc, \Tidentifier(\name), \punnode{\Nparamsopt}, \\
   \wrappedline\ \punnode{\Nfuncargs}, \punnode{\Nrecurselimit}, \punnode{\Nfuncbody}
   \end{array}
   \right)}{\vparsednode}
  \right)
\astarrow\\
{
\overname{
 \left[\DFunc\left(\left\{
 \begin{array}{rcl}
   \FIELDname &:& \name,\\
   \FIELDparameters &:& \astof{\vparamsopt},\\
   \FIELDargs &:& \astof{\vfuncargs},\\
   \funcbody &:& \astof{\vfuncbody},\\
   \FIELDreturntype &:& \none,\\
   \funcsubprogramtype &:& \STProcedure,\\
   \FIELDrecurselimit &:& \astof{\vrecurselimit}\\
   \funcbuiltin &:& \False\\
   \funcqualifierfield &:& \astof{\vqualifier}\\
   \funcoverride &:& \astof{\voverride}
 \end{array}
 \right\}\right)\right]
}{\vastnode}}}
\end{mathpar}

\begin{mathpar}
\inferrule[accessor]{
  \buildaccessorbody(\vbody) \astarrow \vaccessorpair \\
  {
  \desugaraccessorpair\left(
    \begin{array}{l}
    \voverride,\\
    \name,\\
    \astof{\vparamsopt},\\
    \astof{\vfuncargs},\\
    \vsetterarg,\\
    \astof{\tty},\\
    \vaccessorpair
    \end{array}
    \right) \astarrow \vastnode
  }
}{
  {
  \begin{array}{r}
    \builddecl\left(\overname{\Ndecl\left(
      \begin{array}{l}
        \punnode{\Noverride}, \Taccessor, \Tidentifier(\name), \punnode{\Nparamsopt},  \\
        \wrappedline\ \punnode{\Nfuncargs},
        \Tbeq, \Tidentifier(\vsetterarg), \punnode{\Nasty}, \\
        \wrappedline\ \namednode{\vbody}{\Naccessorbody}
      \end{array}
    \right)}{\vparsednode}\right)
  \\ \astarrow \vastnode
  \end{array}
  }
}
\end{mathpar}

\ASTRuleDef{AccessorBody}
\RenderRelation{build_accessor_body}


\begin{mathpar}
\inferrule{
  \vparsednode = \Ndecl(\Tbegin, \namednode{\vaccessors}{\Naccessors}, \Tend, \Tsemicolon)\\
  \buildaccessors(\vaccessors) \astarrow \vaccessorpair
}{
    \buildaccessorbody(\vparsednode) \astarrow \vaccessorpair
}
\end{mathpar}

\ASTRuleDef{RecurseLimit}
\RenderRelation{build_recurse_limit}


\begin{mathpar}
\inferrule[limit]{}{
  \buildrecurselimit\left(\overname{\Nrecurselimit(\Trecurselimit, \punnode{\Nexpr})}{\vparsednode}\right)
  \astarrow
  \overname{\some{\astof{\vexpr}}}{\vastnode}
}
\end{mathpar}

\begin{mathpar}
\inferrule[no\_limit]{}{
  \buildrecurselimit\left(\overname{\Nrecurselimit(\emptysentence)}{\vparsednode}\right)
  \astarrow
  \overname{\none}{\vastnode}
}
\end{mathpar}

\ASTRuleDef{TypedIdentifier}
\RenderRelation{build_typed_identifier}

\begin{mathpar}
\inferrule{}{
  \buildtypedidentifier(\overname{\Ntypedidentifier(\Tidentifier(\id), \punnode{\Nasty})}{\vparsednode}) \astarrow \overname{(\id,\astof{\vasty})}{\vastnode}
}
\end{mathpar}

\ASTRuleDef{OptTypedIdentifier}
\RenderRelation{build_opt_typed_identifier}

\begin{mathpar}
\inferrule{
  \vparsednode = \Ntypedidentifier(\Tidentifier(\id), \namednode{\vastyopt}{\option{\Nasty}})\\
  \buildoption[\Nasty](\vastyopt) \astarrow \astversion{\vastyopt}
}{
  {
  \begin{array}{r}
  \buildopttypedidentifier(\vparsednode) \astarrow \\
  \overname{(\id, \astversion{\vastyopt})}{\vastnode}
  \end{array}
  }
}
\end{mathpar}

\ASTRuleDef{ReturnType}
\RenderRelation{build_return_type}


\begin{mathpar}
\inferrule{}{
  \buildreturntype(\overname{\Nreturntype(\Tarrow, \punnode{\Nty})}{\vparsednode}) \astarrow
  \overname{\astof{\tty}}{\vastnode}
}
\end{mathpar}

\ASTRuleDef{ParamsOpt}
\RenderRelation{build_params_opt}


\begin{mathpar}
\inferrule[empty]{}{
  \buildparamsopt(\overname{\Nparamsopt(\epsilonnode)}{\vparsednode}) \astarrow
  \overname{\emptylist}{\vastnode}
}
\end{mathpar}

\begin{mathpar}
\inferrule[non\_empty]{
  \vparsednode = \Nparamsopt(\Tlbrace, \namednode{\ids}{\ClistZero{\Nopttypedidentifier}}, \Trbrace)\\
  \buildclist[\Nopttypedidentifier](\ids) \astarrow \astversion{\ids}
}{
  \buildparamsopt(\vparsednode) \astarrow
  \overname{\astversion{\ids}}{\vastnode}
}
\end{mathpar}

\ASTRuleDef{FuncArgs}
\RenderRelation{build_func_args}


\begin{mathpar}
\inferrule{
  \buildclist[\Ntypedidentifier](\ids) \astarrow \astversion{\ids}
}{
  \buildfuncargs(\overname{\Nfuncargs(\Tlpar, \namednode{\ids}{\ClistZero{\Ntypedidentifier}}, \Trpar)}{\vparsednode}) \astarrow
  \overname{\astversion{\ids}}{\vastnode}
}
\end{mathpar}

\ASTRuleDef{FuncBody}
\RenderRelation{build_func_body}


\begin{mathpar}
\inferrule{}{
  {
  \begin{array}{r}
  \buildfuncbody(\overname{\Nfuncbody(\Tbegin, \namednode{\vstmts}{\Nstmtlist}, \Tend, \Tsemicolon)}{\vparsednode}) \astarrow \\
  \overname{\astof{\vstmtlist}}{\vastnode}
  \end{array}
  }
}
\end{mathpar}

\ASTRuleDef{Accessors}
\RenderRelation{build_accessors}


\begin{mathpar}
\inferrule{
  \buildstmt(\vgetter) \astarrow \astversion{\vgetter} \hva\\\\
  \buildstmt(\vsetter) \astarrow \astversion{\vsetter} \hva\\\\
  {
  \vaccessorpair \eqdef
  \left\{
  \begin{array}{rcl}
    \accessorpairisreadonly &:& \astof{\visreadonly},\\
    \accessorpairgetter    &:& \astversion{\vgetter},\\
    \accessorpairsetter    &:& \astversion{\vsetter}
  \end{array}
  \right\}
  }
}{
  {
  \begin{array}{r}
    \buildaccessors(
      \punnode{\Nisreadonly}, \Tgetter, \namednode{\vgetter}{\Nstmtlist}, \Tend, \Tsemicolon, \\
      \Tsetter, \namednode{\vsetter}{\Nstmtlist}, \Tend, \Tsemicolon
    ) \\
    \astarrow \vaccessorpair
  \end{array}
  }
}
\end{mathpar}

\begin{mathpar}
\inferrule{
  \buildstmt(\vgetter) \astarrow \astversion{\vgetter} \hva\\\\
  \buildstmt(\vsetter) \astarrow \astversion{\vsetter} \hva\\\\
  {
  \vaccessorpair \eqdef
  \left\{
  \begin{array}{rcl}
    \accessorpairisreadonly &:& \astof{\visreadonly},\\
    \accessorpairgetter    &:& \astversion{\vgetter},\\
    \accessorpairsetter    &:& \astversion{\vsetter}
  \end{array}
  \right\}
  }
}{
  {
  \begin{array}{r}
    \buildaccessors(
      \Tsetter, \namednode{\vsetter}{\Nstmtlist}, \Tend, \Tsemicolon, \\
      \punnode{\Nisreadonly}, \Tgetter, \namednode{\vgetter}{\Nstmtlist}, \Tend, \Tsemicolon
    ) \\
    \astarrow \vaccessorpair
  \end{array}
  }
}
\end{mathpar}

\ASTRuleDef{DesugarAccessorPair}
\RenderRelation{desugar_accessor_pair}


\begin{mathpar}
\inferrule{
  \vqualifier \eqdef \ifthenelseop[H]{\vaccessorpair.\accessorpairisreadonly}{\some{\Readonly}}{\none} \hva\\\\
  {
  \vgetter \eqdef
      \DFunc\left\{
      \begin{array}{rcl}
        \FIELDname &:& \name,\\
        \FIELDparameters &:& \vparameters,\\
        \FIELDargs &:& \vargs,\\
        \funcbody &:& \vaccessorpair.\accessorpairgetter,\\
        \FIELDreturntype &:& \some{\tty},\\
        \funcsubprogramtype &:& \STGetter,\\
        \FIELDrecurselimit   &:& \none\\
        \funcbuiltin &:& \False\\
        \funcqualifierfield &:& \vqualifier\\
        \funcoverride &:& \voverride
      \end{array}
      \right\}
  } \hva\\\\
  \vsetterargs \eqdef [(\vsetterarg, \tty)] \concat \vargs \hva\\\\
  {
  \vsetter \eqdef
      \DFunc\left\{
      \begin{array}{rcl}
        \FIELDname &:& \name,\\
        \FIELDparameters &:& \vparameters,\\
        \FIELDargs &:& \vsetterargs,\\
        \funcbody &:& \vaccessorpair.\accessorpairsetter,\\
        \FIELDreturntype &:& \none,\\
        \funcsubprogramtype &:& \STSetter,\\
        \FIELDrecurselimit   &:& \none\\
        \funcbuiltin &:& \False\\
        \funcqualifierfield &:& \none\\
        \funcoverride &:& \voverride
      \end{array}
      \right\}
  }
}{
  {
    \desugaraccessorpair\left(
      \begin{array}{l}
        \voverride, \\
        \name, \\
        \vparameters, \\
        \vargs, \\
        \vsetterarg, \\
        \tty, \\
        \vaccessorpair
      \end{array}
    \right)
    \astarrow [\vgetter, \vsetter]
  }
}
\end{mathpar}

\ASTRuleDef{Qualifier}
\RenderRelation{build_func_qualifier}


\begin{mathpar}
\inferrule[none]{}{
  \buildqualifier(\overname{\Nqualifier(\emptysentence)}{\vparsednode}) \astarrow \none
}
\end{mathpar}

\begin{mathpar}
\inferrule[pure]{}{
  \buildqualifier(\overname{\Nqualifier(\Tpure)}{\vparsednode}) \astarrow \some{\Pure}
}
\end{mathpar}

\begin{mathpar}
\inferrule[readonly]{}{
  \buildqualifier(\overname{\Nqualifier(\Treadonly)}{\vparsednode}) \astarrow \some{\Readonly}
}
\end{mathpar}

\begin{mathpar}
\inferrule[noreturn]{}{
  \buildqualifier(\overname{\Nqualifier(\Tnoreturn)}{\vparsednode}) \astarrow \some{\Noreturn}
}
\end{mathpar}

\ASTRuleDef{IsReadonly}
\RenderRelation{build_is_readonly}


\begin{mathpar}
\inferrule[none]{}{
  \buildisreadonly(\overname{\Nisreadonly(\emptysentence)}{\vparsednode}) \astarrow \False
}
\end{mathpar}

\begin{mathpar}
\inferrule[readonly]{}{
  \buildqualifier(\overname{\Nisreadonly(\Treadonly)}{\vparsednode}) \astarrow \True
}
\end{mathpar}

\ASTRuleDef{Override}
\RenderRelation{build_override}


\begin{mathpar}
\inferrule[none]{}{
  \buildoverride(\overname{\Noverride(\emptysentence)}{\vparsednode}) \astarrow \none
}
\end{mathpar}

\begin{mathpar}
\inferrule[impdef]{}{
  \buildoverride(\overname{\Noverride(\Timpdef)}{\vparsednode}) \astarrow \some{\Impdef}
}
\end{mathpar}

\begin{mathpar}
\inferrule[implementation]{}{
  \buildoverride(\overname{\Noverride(\Timplementation)}{\vparsednode}) \astarrow \some{\Implementation}
}
\end{mathpar}

%%%%%%%%%%%%%%%%%%%%%%%%%%%%%%%%%%%%%%%%%%%%%%%%%%%%%%%%%%%%%%%%%%%%%%%%%%%%%%%%
\TypeRulesDef{Subprogram Declarations}
%%%%%%%%%%%%%%%%%%%%%%%%%%%%%%%%%%%%%%%%%%%%%%%%%%%%%%%%%%%%%%%%%%%%%%%%%%%%%%%%

The main rule for annotating subprogram declarations is \\
\TypingRuleRef{AnnotateSubprogram}.

We also define the following helper rules:
\begin{itemize}
  \item \TypingRuleRef{AnnotateAndDeclareFunc}
  \item \TypingRuleRef{AnnotateFuncSig}
  \item \TypingRuleRef{AnnotateParams}
  \item \TypingRuleRef{AnnotateOneParam}
  \item \TypingRuleRef{CheckParamDecls}
  \item \TypingRuleRef{FuncSigTypes}
  \item \TypingRuleRef{ParametersOfTy}
  \item \TypingRuleRef{ParametersOfExpr}
  \item \TypingRuleRef{ParametersOfConstraint}
  \item \TypingRuleRef{AnnotateArgs}
  \item \TypingRuleRef{AnnotateOneArg}
  \item \TypingRuleRef{AnnotateReturnType}
  \item \TypingRuleRef{CheckSubprogramPurity}
  \item \TypingRuleRef{DeclareOneFunc}
  \item \TypingRuleRef{SubprogramClash}
  \item \TypingRuleRef{AddNewFunc}
  \item \TypingRuleRef{AddSubprogram}
  \item \TypingRuleRef{Subprogram}
  \item \TypingRuleRef{CheckControlFlow}
  \item \TypingRuleRef{ApproxStmt}
  \item \TypingRuleRef{AllowedAbsConfigs}
\end{itemize}

\TypingRuleDef{AnnotateSubprogram}
See \ExampleRef{Annotating Subprogram Signatures} and
\ExampleRef{Updating the Static Environment for a Subprogram}.


\RenderProseAndFormally{typecheck_decl_func}

\TypingRuleDef{AnnotateAndDeclareFunc}
\RenderRelation{annotate_and_declare_func}

See \ExampleRef{Annotating Subprogram Signatures} and
\ExampleRef{Updating the Static Environment for a Subprogram}.


\RenderProseAndFormally{annotate_and_declare_func}
\CodeSubsection{\AnnotateAndDeclareFuncBegin}{\AnnotateAndDeclareFuncEnd}{../Typing.ml}

\TypingRuleDef{AnnotateFuncSig}
\RenderRelation{annotate_func_sig}

\ExampleDef{Annotating Subprogram Signatures}
\ExampleRef{Annotating Parameters},
\ExampleRef{Checking Parameter Declarations},
\ExampleRef{Annotating Subprogram Arguments},
and
\ExampleRef{Annotating Subprogram Return Types}
show well-typed subprogram signatures and ill-typed subprogram signatures.

The specification in \listingref{AnnotateFuncSig} shows a subprogram
with an associated recursion limit.
\ASLListing{Annotating a subprogram signature with a recursion limit}{AnnotateFuncSig}{\typingtests/TypingRule.AnnotateFuncSig.asl}

The specification in \listingref{AnnotateFuncSig-bad} is ill-typed, since the recursion limit expression \verb|W|
is not a constrained expression.
\ASLListing{An ill-typed subprogram signature}{AnnotateFuncSig-bad}{\typingtests/TypingRule.AnnotateFuncSig.bad.asl}


\RenderProseAndFormally{annotate_func_sig}
\CodeSubsection{\AnnotateFuncSigBegin}{\AnnotateFuncSigEnd}{../Typing.ml}

\TypingRuleDef{AnnotateParams}
\RenderRelation{annotate_params}

\ExampleDef{Annotating Parameters}
In \listingref{AnnotateFuncSig}, the list of explicitly defined parameters
of the function \\
\verb|signature_example| is $\{\texttt{A},\texttt{B}\}$.
Therefore, $\tenvwithparams$ effectively reflects the added declarations \\
\verb|let A: integer{A}| and \verb|let B: integer{B}|.


\RenderProseAndFormally{annotate_params}

\TypingRuleDef{AnnotateOneParam}
\RenderRelation{annotate_one_param}

\ExampleDef{Annotating Subprogram Parameters}
In \listingref{AnnotateOneParam},
the parameters \verb|A| and \verb|B| are annotated as \parameterizedintegertypesterm,
and \verb|C| is annotated as a \wellconstrainedintegertypeterm.
The specification demonstrates how \verb|A|, \verb|B|, and \verb|C| exist as local storage elements
of \verb|parameterized|.
\ASLListing{Annotating subprogram parameters}{AnnotateOneParam}{\typingtests/TypingRule.AnnotateOneParam.asl}

In \listingref{AnnotateOneParam-bad1}, the declaration of the parameter \verb|A| is illegal,
since \verb|A| is also declared as a global storage element.
\ASLListing{An ill-typed subprogram parameter}{AnnotateOneParam-bad1}{\typingtests/TypingRule.AnnotateOneParam.bad1.asl}


\RenderProseAndFormally{annotate_one_param}

\TypingRuleDef{CheckParamDecls}
\RenderRelation{check_param_decls}

\ExampleDef{Checking Parameter Declarations}
The specification in \listingref{CheckParamDecls} is well-typed.
Specifically, the parameter \verb|N|, which must be used in the function signature is used by the
return type.
\ASLListing{A parameter used by the return type}{CheckParamDecls}{\typingtests/TypingRule.CheckParamDecls.asl}

In \listingref{ParametersOfExpr}, the list of extracted parameters \verb|B, C, D, E, F, G|
is equal to the list of declared parameters.

The specification in \listingref{CheckParamDecls-bad} is ill-typed, since the list of extracted parameters is
\verb|D, A, B, C|, while the list of declared parameters is \verb|A, B, C, D|.
\ASLListing{Incorrectly declaring subprogram parameters}{CheckParamDecls-bad}{\typingtests/TypingRule.CheckParamDecls.bad.asl}


\RenderProseAndFormally{check_param_decls}

\TypingRuleDef{ExtractParameters}
\RenderRelation{extract_parameters}

\ExampleDef{The Parameters of a Function}
In \listingref{AnnotateFuncSig}, the set of identifiers that may correspond
to parameters of the function \texttt{signature\_example} is $\{\texttt{A}, \texttt{B}\}$,
since they appear in the type \texttt{bits(A)}
of the argument \texttt{bv} and the type \texttt{bits(A+B)} of the argument \texttt{bv3}.

Finding parameters for each type in the signature of the function \texttt{signature\_example}
yields the following results:
\begin{center}
\begin{tabular}{lll}
\textbf{Expression} & \textbf{Result} & \textbf{Reason}\\
\hline
\texttt{bits(A)} & $\{\texttt{A}\}$ & \texttt{A} is a variable expression \\
& & and \texttt{A} is not defined in the environment.\\
\texttt{bits(W)} & $\emptyset$ & \texttt{W} is defined in the environment.\\
\texttt{bits(A+B)} & $\{\texttt{A}, \texttt{B}\}$ & \texttt{A} and \texttt{B} are variables \\
& & and neither is defined in the environment.\\
\end{tabular}
\end{center}

\ExampleDef{Extracting Parameters from Subprogram Signatures}
The list of parameters extracted from the signature of \verb|parameters_of_expressions|
in \listingref{ParametersOfExpr} is
\verb|B, C, D, E, F, G|.

\ExampleDef{Unsupported Expressions in Subprogram Signatures}
The specification in \listingref{ExtractParameters-bad1} is ill-typed, since the
expressions \\
\verb|N as integer{8,16,32}| and \verb|(N + 1) as integer{9,17}|
are unsupported inside\\ \bitvectortypesterm{} of arguments.
\ASLListing{An ill-typed subprogram signature}{ExtractParameters-bad1}{\typingtests/TypingRule.ExtractParameters-bad1.asl}


\RenderProseAndFormally{extract_parameters}

\TypingRuleDef{FuncSigTypes}
\RenderRelation{func_sig_types}

\ExampleDef{Listing Signature Types}
\listingref{FuncSigTypes} shows the list of signature types for a procedure and a function
in comments above their declarations.
\ASLListing{Listing signature types}{FuncSigTypes}{\typingtests/TypingRule.FuncSigTypes.asl}


\RenderProseAndFormally{func_sig_types}

\TypingRuleDef{ParametersOfTy}
\RenderRelation{paramsofty}

\ExampleDef{The Parameters Extracted from Argument Types}
The parameters extracted from the types of arguments appearing in the signature of \verb|parameters_of_types|
(starting with the return type)
in \listingref{ParametersOfTy} are as follows:
\begin{center}
\begin{tabular}{ll}
\textbf{Type} & \textbf{Parameters}\\
\hline
\verb|bits(A)|            & $[\texttt{A}]$\\
\verb|(bits(B), bits(C))| & $[\texttt{B}, \texttt{C}]$\\
\verb|integer{D..E}|      & $[\texttt{D}, \texttt{E}]$\\
\verb|real|               & $\emptylist$\\
\verb|integer|            & $\emptylist$\\
\end{tabular}
\end{center}

\ASLListing{The parameters extracted from argument types}{ParametersOfTy}{\typingtests/TypingRule.ParametersOfTy.asl}

\ExampleRef{Ill-typed Type Declarations} shows an ill-typed argument type.


\RenderProseAndFormally{paramsofty}

\TypingRuleDef{ParametersOfExpr}
\RenderRelation{params_of_expr}

\ExampleDef{Extracting the Parameters from Expressions}
In \listingref{ParametersOfExpr}, extracting parameters from expressions appearing
in the types of signature arguments (starting with the return type) yield the following lists of parameters:
\begin{center}
\begin{tabular}{ll}
\textbf{Expression} & \textbf{Parameters}\\
\hline
\verb|A|      & $\emptylist$\\
\verb|B|      & $[\texttt{B}]$\\
\verb|C|      & $[\texttt{C}]$\\
\verb|-D|      & $[\texttt{D}]$\\
\verb|E+F|      & $[\texttt{E}, \texttt{F}]$\\
\verb|(G)|      & $[\texttt{G}]$\\
\verb|if H == 0 then I else J|      & $[\texttt{H}, \texttt{I}, \texttt{J}]$\\
\end{tabular}
\end{center}
Notice that, since \verb|A| is declared as a global storage element, it is not extracted as a parameter
for the expression \verb|A| appearing in the return type.

\ASLListing{Extracting the parameters from expressions}{ParametersOfExpr}{\typingtests/TypingRule.ParametersOfExpr.asl}

\listingref{ParametersOfExpr-bad} shows examples of expressions that are illegal in types of arguments.
\ASLListing{Ill-typed expressions in signature types}{ParametersOfExpr-bad}{\typingtests/TypingRule.ParametersOfExpr.bad.asl}


\RenderProseAndFormally{params_of_expr}

\TypingRuleDef{ParametersOfConstraint}
\RenderRelation{params_of_constraint}

\ExampleDef{The Parameters Extracted from a Constraint}
In \listingref{ParametersOfTy}, the list of parameters extracted from the constraint
\verb|{D..E}| is $[\texttt{D}, \texttt{E}]$.

See also \ExampleRef{Extracting the Parameters from Expressions}.


\RenderProseAndFormally{params_of_constraint}

\TypingRuleDef{AnnotateArgs}
\RenderRelation{annotate_args}

\ExampleDef{Annotating Subprogram Arguments}
In \listingref{AnnotateFuncSig}, the annotated arguments are
\texttt{bv}, \texttt{bv2}, \texttt{bv3}, and \texttt{C}.


\RenderProseAndFormally{annotate_args}

\TypingRuleDef{AnnotateOneArg}
\RenderRelation{annotate_one_arg}

\ExampleDef{Annotating a Subprogram Argument}
The specification in \listingref{AnnotateOneArg} shows well-typed argument declarations
and demonstrates how they exist as local storage elements in the subprogram body.
\ASLListing{Well-typed subprogram arguments}{AnnotateOneArg}{\typingtests/TypingRule.AnnotateOneArg.asl}

The specification in \listingref{AnnotateOneArg-bad1} is illegal, since arguments cannot share a name
with global storage elements, \verb|b| in this example.
\ASLListing{An ill-typed subprogram argument}{AnnotateOneArg-bad1}{\typingtests/TypingRule.AnnotateOneArg.bad1.asl}

The specification in \listingref{AnnotateOneArg-bad2} is illegal, since arguments cannot be typed
as \collectiontypesterm.
\ASLListing{An ill-typed subprogram argument}{AnnotateOneArg-bad2}{\typingtests/TypingRule.AnnotateOneArg.bad2.asl}


\RenderProseAndFormally{annotate_one_arg}

\TypingRuleDef{AnnotateReturnType}
\RenderRelation{annotate_return_type}

\ExampleDef{Annotating Subprogram Return Types}
In \listingref{MatchFuncRes}, the subprogram \verb|proc| has no return type.
Annotating its return type does not modify the \staticenvironmentterm{}, and the optional annotated return type is $\none$.
In contrast, the subprogram \verb|returns_value| updates the \staticenvironmentterm{} by
binding $\FIELDreturntype$ to $\unconstrainedinteger$ and returning $\unconstrainedinteger$
as the optional return type.

The specification in \listingref{AnnotateReturnType-bad} is ill-typed, since \collectiontypesterm{} are not allowed
as return types.
\ASLListing{An ill-typed return type}{AnnotateReturnType-bad}{\typingtests/TypingRule.AnnotateReturnType.bad.asl}


\RenderProseAndFormally{annotate_return_type}

\TypingRuleDef{CheckSubprogramPurity}
\RenderRelation{check_subprogram_purity}

\ExampleDef{Checking subprogram purity}
In \listingref{SESForSubprogram}, the bodies of each subprogram are checked as follows:
\begin{itemize}
  \item the side effects of the body of \verb|foo| are all \pureterm{};
  \item the side effects of the body of \verb|bar| are all \readonlyterm{};
  \item the side effects of the body of \verb|goo| are not checked; and
  \item the side effects of the body of \verb|baz| are not checked.
\end{itemize}


\RenderProseAndFormally{check_subprogram_purity}

\TypingRuleDef{DeclareOneFunc}
\RenderRelation{declare_one_func}

See \ExampleRef{Updating the Static Environment for a Subprogram}.


\RenderProseAndFormally{declare_one_func}
\CodeSubsection{\DeclareOneFuncBegin}{\DeclareOneFuncEnd}{../Typing.ml}


\TypingRuleDef{SubprogramClash}
\RenderRelation{subprogram_clash}

The function assumes there exists a binding for $\name$ in the
$\subprograms$ map of $\tenv$.

\ExampleDef{Subprogram Clashing}
In \listingref{SubprogramTypesClash}, the getter for \verb|X| and the setter for \verb|X| do not clash,
since they have non-clashing subprogram types (see \TypingRuleRef{SubprogramTypesClash}).

In \listingref{SubprogramTypesClash-bad1}, the getter for \verb|X| and the function \verb|X| clash,
since their subprogram types clash and their signatures clash (see \TypingRuleRef{HasArgClash}).

In \listingref{SubprogramTypesClash-bad2}, the function \verb|X| and the procedure \verb|X| clash
since their subprogram types clash and their signatures clash.

In \listingref{SubprogramClash-bad1}, the function \verb|X| and the procedure \verb|X| clash since subprogram types clash and their qualifiers differ.
\ASLListing{Clashing subprogram qualifiers}{SubprogramClash-bad1}{\typingtests/TypingRule.SubprogramClash.bad1.asl}


\RenderProseAndFormally{subprogram_clash}


\TypingRuleDef{SubprogramTypesClash}
\RenderRelation{subprogram_types_clash}

\ExampleDef{Clashing Subprogram Types}

The specification in \listingref{SubprogramTypesClash} contains an accessor for \verb|X|.
The subprogram type of its getter does not clash with the subprogram type of its setter.
\ASLListing{Non-clashing subprogram types}{SubprogramTypesClash}{\typingtests/TypingRule.SubprogramTypesClash.asl}

The specification in \listingref{SubprogramTypesClash-bad1} contains an accessor for \verb|X|.
The subprogram types of its getter and setter both clash with the subprogram type of the procedure\\ \verb|X| ($\STProcedure$).
\ASLListing{Clashing subprogram types}{SubprogramTypesClash-bad1}{\typingtests/TypingRule.SubprogramTypesClash.bad1.asl}

The specification in \listingref{SubprogramTypesClash-bad2} contains a function \verb|X| and a procedure \verb|X|
whose subprogram types ($\STFunction$ and $\STProcedure$, respectively) clash.
\ASLListing{Clashing subprogram types}{SubprogramTypesClash-bad2}{\typingtests/TypingRule.SubprogramTypesClash.bad2.asl}

\ProseParagraph
\ProseEqdef{$\vb$}{$\True$ unless
$\vsone$ is $\STGetter$ and $\vstwo$ is $\STSetter$ or
$\vsone$ is $\STSetter$ and $\vstwo$ is $\STGetter$}.

\FormallyParagraph
\begin{mathpar}
\inferrule{
  {
  \vb \eqdef (\vsone, \vstwo) \not\in
  \left\{
  \begin{array}{l}
    (\STGetter,       \STSetter),\\
    (\STSetter,       \STGetter)
  \end{array}
  \right\}
  }
}{
  \subprogramtypesclash(\subpgmtypeone, \subpgmtypetwo) \typearrow \vb
}
\end{mathpar}

\TypingRuleDef{AddNewFunc}
\RenderRelation{add_new_func}

The function $\addnewfunc$ updates the $\overloadedsubprograms$ map in the \staticenvironmentterm{},
and must be followed by $\addsubprogram$, which updates the $\subprograms$ map (see \TypingRuleRef{AddSubprogram}).

\ExampleDef{Updating the Static Environment for a Subprogram}
The specification in \listingref{DeclareSubprograms}
contains two subprograms named \verb|add_10| with the signatures
\verb|func add_10(x: integer) => integer| and
\verb|func add_10(x: real) => real|, which for brevity we will refer to by
$I$ and $R$, respectively.
These subprograms are added to the \staticenvironmentterm{}
based on the dependency ordering \\
(see \TypingRuleRef{DeclDependencies}).
In this case, there is no ordering between these two subprograms.
Assuming $I$ is processed before $R$, adding $I$ to the
\staticenvironmentterm{}, the \staticenvironmentterm{} has
$\addten \not\in \dom(\overloadedsubprograms)$
and is then updated by binding $\addten$ to $\{\addten\}$ in $\overloadedsubprograms$.
Then, processing $R$, results in
binding $\addten$ to
$\{\addten, \addtenone\}$ where \addtenone{} will be used for $R$.

\ExampleDef{Illegal Subprogram Declarations due to Clashes}
The specifications in \listingref{AddNewFunc-bad1},
\listingref{AddNewFunc-bad2}, and
\listingref{AddNewFunc-bad3}
are ill-typed due to subprogram declaration clashes.
\ASLListing{Subprogram clashes}{AddNewFunc-bad1}{\typingtests/TypingRule.AddNewFunc.bad1.asl}
\ASLListing{Subprogram clashes}{AddNewFunc-bad2}{\typingtests/TypingRule.AddNewFunc.bad2.asl}
\ASLListing{Subprogram clashes}{AddNewFunc-bad3}{\typingtests/TypingRule.AddNewFunc.bad3.asl}


\RenderProseAndFormally{add_new_func}

\CodeSubsection{\AddNewFuncBegin}{\AddNewFuncEnd}{../Typing.ml}

\TypingRuleDef{Subprogram}
\RenderRelation{annotate_subprogram}

Note that the return type in $\vf$ has already been annotated by $\annotatefuncsig$.

\ExampleDef{Annotating Subprograms}
\listingref{typing-subprogram} shows an example of a well-typed procedure --- \verb|my_procedure|
and an example of a well-typed function --- \verb|flip_bits|.
\ASLListing{Typing subprograms}{typing-subprogram}{\typingtests/TypingRule.Subprogram.asl}


\RenderProseAndFormally{annotate_subprogram}


\CodeSubsection{\SubprogramBegin}{\SubprogramEnd}{../Typing.ml}

\TypingRuleDef{CheckControlFlow}
\RenderRelation{check_control_flow}

\ExampleDef{Ensuring All Terminating Paths Terminate Correctly}
In \listingref{CheckControlFlow}, every evaluation of the function body for\\
\verb|all_terminating_paths_correct| terminates by either returning
a value, throwing an exception, or evaluating an \unreachablestatementterm.
\ASLListing{All terminating paths terminate correctly}{CheckControlFlow}{\typingtests/TypingRule.CheckControlFlow.asl}

In \listingref{CheckControlFlow-bad2}, the path through the function body
for\\ \verb|incorrect_terminating_path|
where \verb|v != Zeros{N}| evaluates to
$\True$ and \verb|flag| evaluates to $\False$ terminates without
returning a value, throwing an exception, or evaluating an \unreachablestatementterm,
which is a \typingerrorterm.
\ASLListing{An incorrectly terminating path}{CheckControlFlow-bad2}{\typingtests/TypingRule.CheckControlFlow.bad2.asl}

\ExampleDef{Illegally returning from a \texttt{noreturn} Subprogram}
The specification in \listingref{CheckControlFlow-bad3}
is ill-typed as \verb|returning| is declared with the \verb|noreturn|
qualifier, yet it returns via a \passstatementterm.
\ASLListing{Illegally returning from a \texttt{noreturn} subprogram}{CheckControlFlow-bad3}{\typingtests/TypingRule.CheckControlFlow.bad3.asl}


\RenderProseAndFormally{check_control_flow}

\TypingRuleDef{AllowedAbsConfigs}
\RenderRelation{allowed_abs_configs}

\ExampleDef{The Allowed Abstract Configurations for a Subprogram}
The set of \Proseabstractconfigurations{} allowed for the subprogram \\
\verb|all_terminating_paths_correct|
in \listingref{CheckControlFlow} is $\{\AbsAbnormal, \AbsReturning\}$, since it is
a function.

The set of \Proseabstractconfigurations{} allowed for the subprogram \verb|returning|
in \listingref{CheckControlFlow-bad3} is $\{\AbsAbnormal\}$, since it is qualified
by \verb|noreturn|.


\RenderProseAndFormally{allowed_abs_configs}

\TypingRuleDef{ApproxStmt}
\RenderRelation{approx_stmt}

The type \RenderType{abstract_configuration}
of \emph{\Proseabstractconfigurations{}} consists of symbols that represent
subsets of the set of dynamic configurations that may be returned by
$\evalstmt((\tenv, \cdot), \vs)$:
\begin{itemize}
  \item $\AbsContinuing$ represents $\TContinuing$;
  \item $\AbsReturning$ represents $\TReturning$; and
  \item $\AbsAbnormal$ represents $\TThrowing \cup \TDynError$.
\end{itemize}

\ExampleDef{Determining the Abstract Configurations of Statements}
In \listingref{ApproxStmt}, the function bodies of all functions
terminate by either returning a value, throwing an exception, or executing
the \unreachablestatementterm.
\ASLListing{Determining the abstract configurations of statements}{ApproxStmt}{\typingtests/TypingRule.ApproxStmt.asl}

In \listingref{ApproxStmt-bad1}, the function body of \verb|loop_forever|
is ill-typed, since the conservative analysis of \TypingRuleRef{ApproxStmt}
cannot determine that the \whilestatementterm{} never terminates.
To make the function body well-typed, another statement following the loop
can be added, for example, an \unreachablestatementterm.
\ASLListing{An ill-typed function body}{ApproxStmt-bad1}{\typingtests/TypingRule.ApproxStmt.bad1.asl}

\RenderProseAndFormally{approx_stmt}
